\documentclass[11pt]{article}
\usepackage{amsmath,amssymb,amsthm,mathtools}
\usepackage[margin=1in]{geometry}
\usepackage{enumitem}
\usepackage{hyperref}

\newcommand{\R}{\mathbb R}
\newcommand{\C}{\mathbb C}
\newcommand{\Z}{\mathbb Z}
\newcommand{\Q}{\mathbb Q}
\newcommand{\HH}{\mathcal H}
\newcommand{\W}{\mathcal W}

\title{Status of the \(S^3\)--Connes Program for the Riemann Hypothesis}
\date{10 August 2026}

\begin{document}
\maketitle

\begin{abstract}
The purpose of this program is not to construct a new explicit formula for
\(\zeta\), but to use an \(S^3\)-based geometry to understand the concrete
positivity obstruction in Connes' spectral approach to the Riemann
Hypothesis (RH).  The main progress is a reduction of the observed
near-null direction in the Connes--Consani--Moscovici (CCM) finite
approximations to a very specific boundary/tunnelling phenomenon in the
prolate decomposition.  For the distinguished CCM test vector
\(h_\lambda\), two potentially dangerous scalar defects are now controlled
at precisely the prolate tunnelling scale:
\[
        \widehat h_\lambda(0)=0
        \qquad\text{exactly},
        \qquad
        h_\lambda(0)=O(1-\chi_4).
\]
Thus the remaining issue is no longer a vague approximation problem.
The principal unresolved estimate is
\[
        QW_\lambda(Eh_\lambda)\asymp 1-\chi_4,
\]
or a sufficiently strong one-sided version of it.  Here \(E\) is the
arithmetic periodization operator, \(W_\lambda\) is the truncated Weil
operator/form, \(Q\) removes the expected low-dimensional exceptional
sector, and \(\chi_4\) is the relevant prolate concentration eigenvalue.
Extracting this estimate from the 2023 trace formula is currently the
shortest plausible route to turning the geometric mechanism into a
rigorous statement directly relevant to RH.
\end{abstract}

\section{The Connes problem}

We use Weil's positivity criterion as the analytic target.  In the
multiplicative picture the arithmetic operator
\[
        E(f)(u)
        =
        u^{1/2}\sum_{n\geq1}f(nu)
\]
has Mellin transform
\[
        \widehat{Ef}(s)=\zeta(s)\widehat f(s),
\]
up to the chosen normalization.  Equivalently,
\[
        E(f)(u)
        =
        u^{1/2}\operatorname{Tr}_{\ell^2(\mathbb N)}
        f\!\left(ue^{H_{\rm BC}}\right),
\]
so the zeta factor is built directly into the spectral/arithmetic
periodization.

Connes' program rewrites the explicit formula as a trace/positivity
problem.  Schematically,
\[
        \mathrm{RH}
        \quad\Longleftrightarrow\quad
        W(F)\geq0
\]
on the appropriate codimension-two space of test functions, the two
removed directions corresponding to the familiar pole/end-point
obstructions.

The difficult point is therefore not finding the zeros of \(\zeta\):
it is proving positivity after arithmetic periodization and removal of
the exceptional modes.

CCM's finite-\(\lambda\) calculations reveal a striking almost-null
direction.  Their numerical Figure~4 suggests that its defect is
controlled by a prolate concentration eigenvalue.  Explaining this
phenomenon rigorously is the immediate target of the present program.

\section{Where the \(S^3\) geometry enters}

The \(S^3\) construction supplies a geometric model for the same
decomposition that appears analytically in the Connes--CCM problem.
The useful dictionary is schematic rather than an asserted literal
equivalence:
\[
\begin{array}{c|c}
S^3\text{ geometry} & \text{Connes--CCM analysis}\\
\hline
\text{global compact geometry}
&
\text{completed spectral problem}
\\
\text{two distinguished/boundary directions}
&
\text{two exceptional Weil directions}
\\
\text{geometric concentration}
&
\text{prolate concentration}
\\
\text{failure to remain in a geometric sector}
&
\text{time--band leakage}
\\
\text{small boundary/tunnelling term}
&
1-\chi_n
\\
\text{geometric projection}
&
Q\text{ / removal of exceptional modes}.
\end{array}
\]

The important conceptual prediction from the geometry is therefore:
\emph{after the exact exceptional directions have been removed, the
remaining failure of positivity should be a boundary leakage phenomenon,
not a bulk arithmetic phenomenon.}

This is now substantially supported by the analytic reductions below.

\section{The distinguished CCM vector}

Let \(h_{0,\lambda}\) and \(h_{4,\lambda}\) denote the two relevant
prolate modes and set
\[
        h_\lambda
        =
        \alpha h_{0,\lambda}
        +
        \beta h_{4,\lambda},
\]
with the CCM choice of coefficients.  Let
\(\chi_0,\chi_4\) be their concentration eigenvalues, with
\(\chi_4\to1\) as \(\lambda\to\infty\).

The arithmetic candidate is
\[
        k_\lambda=E(h_\lambda).
\]
This is the finite-scale vector numerically behaving like the ground
state of the simple-even Weil form.

The key question had been whether the exceptional components generated
by \(E\) were much larger than the exponentially/super-polynomially
small prolate leakage.  Direct norm estimates on \(E\) do not resolve
this: crude leakage estimates are only polynomial and therefore lose
the scale visible numerically.

The geometry suggested looking instead for exact cancellation of the
exceptional components.

\section{New reductions}

For the CCM combination one obtains
\[
        \boxed{\widehat h_\lambda(0)=0.}
\]
Thus one exceptional scalar direction vanishes \emph{exactly}, rather
than merely being small.

For the second endpoint one finds
\[
 h_\lambda(0)
 =
 \alpha h_{0,\lambda}(0)
 \frac{\chi_4-\chi_0}{\chi_4},
\]
and hence, since the relevant concentration eigenvalues approach one
together,
\[
        \boxed{h_\lambda(0)=O(1-\chi_4).}
\]

Consequently the two obvious exceptional defects satisfy
\[
\boxed{
\text{exceptional sector}
=
0+O(1-\chi_4).
}
\]

This is the main new structural result.  It converts the previously
problematic mismatch of scales into exactly the scale predicted by the
\(S^3\)/prolate geometry.

In particular, the numerical near-kernel is no longer mysterious:
the CCM vector has been engineered, or geometrically forced, to kill
one obstruction exactly and suppress the other to the tunnelling scale.

\section{The remaining lemma}

Let \(W_\lambda\) denote the finite/truncated Weil operator or quadratic
form in the CCM setup, and let \(Q\) denote projection away from the
exceptional sector.  The desired statement is now essentially
\[
        \boxed{
        QW_\lambda(Eh_\lambda)
        \asymp
        1-\chi_4.
        }
        \tag{*}
\]
Depending on normalization, a one-sided coercive estimate of the
corresponding quadratic form may suffice.

The significance of \((*)\) is that both sides now live at the same
natural scale.  Earlier approaches attempted to compare a generic
periodization error with \(1-\chi_4\), which is hopelessly small.
The endpoint identities show that the large apparent errors lie in
directions which cancel before the decisive comparison is made.

Thus the problem has been reduced from

\[
\text{``control the arithmetic operator }E\text{ extremely accurately''}
\]

to

\[
\text{``identify the surviving boundary term in the trace formula.''}
\]

That is a much sharper problem.

\section{Relation to the 2023 trace formula}

The next step should be to insert \(h_\lambda\), or the corresponding
rank-one/prolate projector, directly into the 2023 trace formula and
separate
\[
        \text{bulk}
        +
        \text{exceptional terms}
        +
        \text{boundary crossing}.
\]

The first two pieces should cancel or be removed by \(Q\).  The claim
suggested by both the geometry and the CCM numerics is that what remains
is precisely the prolate boundary leakage:
\[
        \text{boundary contribution}
        \sim
        1-\chi_4.
\]

There is already an exact analytic indication in this direction.
For the prolate Sturm--Liouville operator \(L_c\) and the relevant
concentration projector \(P\),
\[
        [L_c,P]=0.
\]
When the prime part of the explicit formula is written in terms of
multiplicative translations, commutator matrix elements localize to
boundary-crossing strips.  Thus the nonlocal arithmetic operation becomes
a boundary phenomenon after comparison with the prolate decomposition.

This is strikingly close to the mechanism predicted by the \(S^3\)
geometry.

\section{What has not worked}

The principal unsuccessful routes can be summarized briefly.

\begin{itemize}[leftmargin=2em]
\item \textbf{Naive \(E\)-norm estimates.}
They give errors far larger than \(1-\chi_4\) and therefore cannot
explain the CCM numerics.

\item \textbf{Generic Fourier/prolate leakage bounds.}
These control the wrong object.  The decisive cancellation occurs in
the exceptional scalar components before generic norm estimates become
relevant.

\item \textbf{Prime--archimedean cancellation at low frequency.}
A direct symbol calculation does not yield the hoped-for cancellation;
the completed prime and \(\Gamma\)-terms can reinforce rather than
cancel there.

\item \textbf{Commutator estimates alone.}
The identity/localization to boundary-crossing strips is useful, but
without the trace-formula coefficient calculation it does not determine
the sign or exact size required for Weil positivity.
\end{itemize}

These failures are informative: they point away from global estimates
and toward an exact boundary-term computation.

\section{Current interpretation}

The emerging picture is
\[
\boxed{
\begin{gathered}
S^3\text{ geometry}
\\
\Downarrow
\\
\text{two exceptional directions + boundary leakage}
\\
\Downarrow
\\
\text{prolate decomposition}
\\
\Downarrow
\\
\widehat h_\lambda(0)=0,\qquad
h_\lambda(0)=O(1-\chi_4)
\\
\Downarrow
\\
\text{only a tunnelling-scale Weil defect remains.}
\end{gathered}}
\]

This gives a geometric explanation for why the CCM construction should
work: the small eigenvalue is not an accidental numerical approximation
to zeta.  It appears to measure the failure of a nearly invariant
geometric/prolate sector to close exactly at finite scale.

In this sense the \(S^3\) construction is currently most valuable not
as a replacement for Connes' machinery, but as a structural guide to
\emph{which term in that machinery must be estimated and why it should
have the correct sign and scale}.

\section{Shortest plausible route forward}

The research program should now be aggressively concentrated on the
following chain.

\begin{enumerate}[leftmargin=2em]
\item \textbf{Extract the boundary term.}
Specialize the 2023 trace formula to the CCM vector \(h_\lambda\) and
write \(QW_\lambda(Eh_\lambda)\) explicitly after using
\[
        \widehat h_\lambda(0)=0,
        \qquad
        h_\lambda(0)=O(1-\chi_4).
\]

\item \textbf{Prove the tunnelling estimate.}
Show
\[
        QW_\lambda(Eh_\lambda)
        =
        C_\lambda(1-\chi_4)
        +o(1-\chi_4),
\]
ideally with
\[
        0<c\leq C_\lambda\leq C.
\]
The sign of \(C_\lambda\) is more important than a sharp asymptotic
constant.

\item \textbf{Identify the coefficient geometrically.}
Relate \(C_\lambda\) to the \(S^3\) boundary/intersection quantity.
This is where the geometry could turn an analytic estimate into a
structural positivity statement rather than a delicate asymptotic
calculation.

\item \textbf{Upgrade the single-vector calculation to coercivity.}
Once the CCM ground-state direction is controlled, determine whether
the orthogonal complement already has a uniform positive gap from the
existing trace-formula/prolate theory.  If so, the entire positivity
problem reduces to the one tunnelling direction.

\item \textbf{Pass to the limit.}
Combine finite-\(\lambda\) positivity with density/convergence of the
CCM approximation to obtain Weil positivity on the full admissible
test-function space, hence RH.
\end{enumerate}

The critical logical point is Step~4.  Proving \((*)\) explains
Figure~4 and controls the apparent ground state, but RH follows only
if the rest of the Weil form is simultaneously known to be positive
with adequate uniform control.  This distinction should be maintained
explicitly.

\section{Research target}

The immediate target can therefore be stated as a compact lemma.

\medskip
\noindent
\textbf{Boundary Leakage Lemma.}
\emph{For the CCM prolate combination \(h_\lambda\), after projection
away from the two exceptional Weil directions, arithmetic
periodization converts the finite Weil defect into a boundary term
comparable to the fourth prolate tunnelling probability:}
\[
        \left\|QW_\lambda(Eh_\lambda)\right\|
        \asymp
        1-\chi_4,
\]
\emph{with the corresponding quadratic contribution having the
positivity sign required by Weil's criterion.}

\medskip

A proof of this lemma would provide a rigorous explanation of the
numerical phenomenon observed by CCM and would establish the first
direct analytic bridge from the \(S^3\) geometry to the concrete
obstruction in Connes' RH program.

The most efficient next calculation is therefore not another global
estimate and not a new numerical experiment.  It is the explicit
specialization of the 2023 trace formula to \(Eh_\lambda\), retaining
the boundary-crossing term and exploiting the two endpoint
cancellations above.  If that computation produces a positive multiple
of \(1-\chi_4\), the program should immediately turn to proving a
uniform gap on its orthogonal complement; together those two statements
would place the approach on a direct route to Weil positivity and RH.

\end{document}
